\section{Stochastic modelling}

Generally, at the outset of an infection, a relatively small number of colony forming units (bacteria or virus) will enter the body. 
Be it during inhalation into the respiratory system, through the skin with an insect's bite, or via ingestion, those initial pathogenic particles will be critical in determining the fate of the infection's success. 

The founding population will either be depleted to extinction, or expand to a capacity where its growth becomes established and requires a more targeted effort of the immune system to deal with. Whether they succeed in surviving beyond the first few hours of infection will depend on causes of a chance nature; encounters with cells and other immune factors that are subject to the randomness of Brownian motion.  


Stochastic effects play a more significant role in a population which is small. Because as the count of individuals increases, the dynamics of their overall growth and decline will become increasingly subject to the law of large numbers, and the expected fate of particular bacteria becomes telling of the group's trajectory. 

For this reason, because we are interested in the early stages of a Yersinia Pestis infection, it will be important to use stochastic models in our analysis.
